Firm characteristics are commonly represented as fixed vectors, even though
evidence about firms' operations arrives as heterogeneous collections of
articles and reports.
We study how distances between distributions of firm characteristics restrict
systematic covariance.
For given latent exposure laws, quadratic Wasserstein geometry yields sharp
covariance endpoints over admissible couplings.
Under a common randomized bi-Lipschitz characteristic-to-exposure map, bounded
risk-coordinate slack, and a maintained return-covariance bridge,
characteristic-side distance yields a conditional interval for the covariance ceiling.
Greater separation then leaves less scope for aligned systematic exposures when
these assumptions are tight, while a nonnegative transport-excess term records
the gap between realized and maximum-covariance arrangements.
The resulting envelope is scenario analysis rather than an expected-return,
no-arbitrage, or identified structural model.
Empirically, we proxy characteristic distributions with encoder-only
language-model embeddings of financial news articles, disclosures, and analyst
reports for Nasdaq firms over
\ppvalue{sample-year-start}--\ppvalue{sample-year-end}.
In a dyadic regression with symmetric firm effects, greater pairwise
article-embedding W$_2$ distance is associated with weaker return co-movement:
the coefficient is
\ppnum[3]{confound-w2-coef}, with 95\% interval
[\ppnum[3]{confound-w2-ci-low}, \ppnum[3]{confound-w2-ci-high}].
The estimate retains the same sign when returns are measured in a later window,
although the feature construction is not point-in-time and the evidence remains
reduced form.
